%%%%%%%%%%%%%%%%%%%%%%%%%%%%%%%%%%%%%%%%%%%%%%%%%%%%%%%%%%%%%%%%%%%%%%%%
%                                                                      %
%     File: Thesis_Appendix_A.tex                                      %
%     Tex Master: Thesis.tex                                           %
%                                                                      %
%     Author: Francisco Castro                                         %
%     Last modified :  29 Jan 2024                                     %
%                                                                      %
%%%%%%%%%%%%%%%%%%%%%%%%%%%%%%%%%%%%%%%%%%%%%%%%%%%%%%%%%%%%%%%%%%%%%%%%

\chapter{Annex A}
\label{chapter:annexA}

\section{Euler angles}
\label{section:eulerAngles}

Euler rotation angles defining an intrinsic $z-y-x$ (yaw-pitch-roll) rotation, with a rotation sequence of $z-y'-x''$ around the moving frame $\BodyFrame$ is denoted as,

\begin{equation}
	\mathbf{\RotMatrix}_{\BodyFrame}^{\WorldFrame} = \bm{\RotMatrix}_z (\yaw) \bm{\RotMatrix}_y (\pitch) \bm{\RotMatrix}_x (\roll)
\end{equation}

, with the previous multiplication to yield a mapping from frame $\BodyFrame$ to $\WorldFrame$, which stands as,

\begin{equation}
	\begin{aligned}
	\mathbf{\RotMatrix}_{\BodyFrame}^{\WorldFrame} (\roll, \pitch, \yaw) &=
	\begin{bmatrix}
		\cos(\yaw) & -\sin(\yaw) & 0 \\
		\sin(\yaw) & \cos(\yaw) & 0 \\
		0 & 0 & 1
	\end{bmatrix}
	\begin{bmatrix}
		\cos(\pitch) & 0 & \sin(\pitch) \\
		0 & 1 & 0 \\
		-\sin(\pitch) & 0 & \cos(\pitch)
	\end{bmatrix}
	\begin{bmatrix}
		1 & 0 & 0 \\
		0 & \cos(\roll) & -\sin(\roll) \\
		0 & \sin(\roll) & \cos(\roll)
	\end{bmatrix} \\
	& =
	\begin{bmatrix} 
		\cos(\yaw) \cos(\pitch) & \cos(\yaw) \sin(\pitch) \sin(\phi) - \sin(\yaw) \cos(\phi) & \cos(\yaw) \sin(\pitch) \cos(\roll) + \sin(\yaw) \sin(\roll) \\
		\sin(\yaw) \cos(\pitch) & \sin(\yaw) \sin(\pitch) \sin(\roll) + \cos(\yaw) \cos(\roll) & \sin(\yaw) \sin(\pitch) \cos(\roll) - \cos(\yaw) \sin(\roll) \\
		-\sin(\pitch) & \cos(\pitch) \sin(\roll) & \cos(\pitch) \cos(\roll)
	\end{bmatrix}
	\end{aligned}
	\label{eq:rotationMatrixW2B}
\end{equation}

Rotation matrices are singular ($\det(\bm{\RotMatrix}_{\BodyFrame}^{\WorldFrame}) = 1$) and symmetric, $(\bm{\RotMatrix}_{\BodyFrame}^{\WorldFrame})^T = (\bm{\RotMatrix}_{\BodyFrame}^{\WorldFrame})^{-1} = \bm{\RotMatrix}_{\WorldFrame}^{\BodyFrame}$. Euler rotation matrices can be derived from 24 different conventions ranging from extrinsic or intrinsic rotation with different rotation orders. More information on this can be found in \cite{Roithmayr2015}.

\section{Skew-symmetric matrix}
\label{section:skewSymmetricMatrix}

A skew-symmetric matrix $\bm{S}$ is a square matrix whose transpose equals its negative, $\bm{S}^T = -\bm{S}$. In this thesis, three by three matrices are employed which take the following form,

\begin{equation}
	\bm{S(r} = (x, \; y, \; z)\bm{)} = \begin{bmatrix}
								0 & -z & y\\
								z & 0 & -x\\
								-y & x & 0	
	\end{bmatrix}
	\label{eq:skewSymmetricMatrix} 
\end{equation}

In case this matrix is multiplied by a column vector $\mathbf{a}$, it is the same as applying the cross product between $\mathbf{r}$ and $\mathbf{a}$. Then, it is stated that,

\begin{equation}
	\mathbf{S(r)} \mathbf{a} = \mathbf{r} \times \mathbf{a}
	\label{eq:skewSymmetricMatrixProperties} 
\end{equation}

\section{Inertia tensor}
\label{section:inertiaTensor}

The inertia tensor of a rigid body is given by,

\begin{equation}
	\mathbf{\InertiaTensor} = \begin{bmatrix}
		\InertiaTensor_{xx} & \InertiaTensor_{xy} & \InertiaTensor_{xz}\\
		\InertiaTensor_{yx} & \InertiaTensor_{yy} & \InertiaTensor_{yz}\\
		\InertiaTensor_{zx} & \InertiaTensor_{zy} & \InertiaTensor_{zz}
	\end{bmatrix}
	\label{eq:annexInertiaTensor}
\end{equation}

, where the products of inertia are equal between each other ($I_{xy} = I_{yx}$, $I_{xz} = I_{zx}$, $I_{yz} = I_{zy}$).

\section{Differentiation of a vector}
\label{section:vectorDiff}

If a vector is expressed with respect to a non-inertial frame $\mathcal{P}$, it is still possible to differentiate that same vector with respect to an inertial frame $\mathcal{Q}$ \cite{Ferdinand2015}.

Let $\mathbf{p} = p_1 \; \mathbf{i} + p_2 \; \mathbf{j} + p_3 \; \mathbf{k}$, with $\mathbf{i}, \; \mathbf{j}, \; \mathbf{k}$ to be the canonical vectors of a non-inertial frame $\mathcal{P}$. Then, the differentiation of $\mathbf{p}$ with respect to $\mathcal{Q}$ is given by,

\begin{equation}
	\derivative{\mathbf{p}}\Bigr|_{\mathcal{Q}} = \dot{p}_1 \; \mathbf{i} + \dot{p}_2 \; \mathbf{j} + \dot{p}_3 \; \mathbf{k} + p_1 \derivative{\mathbf{i}} + p_2 \derivative{\mathbf{j}} + p_3 \derivative{\mathbf{k}}
	\label{eq:vectorDiff}
\end{equation}
 
, where $\derivative{\mathbf{p}}\Bigr|_{\mathcal{P}} = \dot{p}_1 \; \mathbf{i} + \dot{p}_2 \; \mathbf{j} + \dot{p}_3 \; \mathbf{k}$ because the canonical vectors $\mathbf{i}, \; \mathbf{j}, \; \mathbf{k}$ do not change with respect to time if the observer is located at the origin of $\mathcal{P}$. According to the general expression $\derivative{\mathbf{r}} = \bm{\Omega} \times \mathbf{r}$, it comes that $\derivative{\mathbf{i}} = \bm{\Omega} \times \mathbf{i}, \; \derivative{\mathbf{j}} = \bm{\Omega} \times \mathbf{j}, \; \derivative{\mathbf{k}} = \bm{\Omega} \times \mathbf{k}$. Then, equation \ref{eq:vectorDiff} may be rewritten as,

\begin{equation}
	\derivative{\mathbf{p}}\Bigr|_{\mathcal{Q}} = \derivative{\mathbf{p}}\Bigr|_{\mathcal{P}} + \mathbf{\Omega} \times \mathbf{p} 
	\label{eq:vectorDiff2}
\end{equation} 

, where $\bm{\Omega}$ denotes the angular velocity of frame $\mathcal{P}$ with respect to frame $\mathcal{Q}$.

\section{Parallel-axis theorem}
\label{section:parallelAxisTheorem}

\section{Quaternions}
\label{section:quaternions}

\subsection{Rotation matrix}
\label{subsection:rotationMatrix}
	
A quaternion is defined by $\Quaternion = \qw + \qx \mathbf{i} + \qy \mathbf{j} + \qz \mathbf{k}$. The right-hand rule convention is adopted.

The rotation matrix \cite{attitudeRepresentations2001} which maps linear velocities from a non-inertial frame $\mathcal{P}$ to an inertial frame $\mathcal{Q}$ is the following,

\begin{equation}
	\mathbf{\bar{\RotMatrix}}_{\mathcal{P}}^{\mathcal{Q}} = \begin{bmatrix}
		1 - 2 (\qy^2 + \qz^2) & 2 (\qx \qy - \qz \qw) & 2 (\qx \qz + \qy \qw) \\
		2 (\qx \qy + \qz \qw) & 1 - 2(\qx^2 + \qz^2) & 2 (\qy \qz - \qx \qw) \\
		2 (\qx \qz - \qy \qw) & 2 (\qy \qz + \qx \qw) & 1 - 2(\qx^2 + \qy^2)  
	\end{bmatrix}
	\label{eq:quatRotationMatrix}
\end{equation}

The rotation matrix which maps angular velocities to quaternion rates is the following,

\begin{equation}
	\mathbf{\bar{\RotMatrix}}_{\mathcal{P}}^{\mathcal{Q}} = \frac{1}{2} \times \begin{bmatrix}
		-\qx & -\qy & -\qz \\
		\qw & -\qz  & \qy \\
		\qz & \qw & -\qx\\
		-\qy & \qx & \qw  
	\end{bmatrix}
	\label{eq:orientationQuatRotationMatrix}
\end{equation}

\subsection{Error matrix}
\label{subsection:errorMatrix}

The difference between two quaternions is expressed in \ref{eq:errorQuaternion}. More information on the derivation of this result can be found in \cite{attitudeError}.

\begin{equation}
	\delta \Quaternion = \Quaternion \otimes \Quaternion_{ref}^{-1} \begin{bmatrix}
		 \qw \qw^{ref} + \qx \qx^{ref} + \qy \qy^{ref} + \qz \qz^{ref} \\
		 -\qw \qx^{ref} + \qx \qw^{ref} - \qy \qz^{ref} + \qz \qy^{ref} \\
		 -\qw \qy^{ref} + \qx \qz^{ref} + \qy \qw^{ref} - \qz \qx^{ref} \\
		 -\qw \qz^{ref} - \qx \qy^{ref} + \qy \qx^{ref} + \qz \qw^{ref} 
	\end{bmatrix}
	\label{eq:errorQuaternion}
\end{equation}

\section{Gauss-Legrendre quadrature}
\label{section:gaussQuadrature}

A Gauss-Legrendre quadrature is an integral approximation method, which yields an exact result for polynomials of degree $2n -1$ or less.

\begin{equation}
  \int_a^b f(x)\, dx \approx \int_{-1}^{1} f\left(\frac{b-a}{2} \xi + \frac{a+b}{2} \right) \frac{dx}{d\xi} d\xi
\label{eq:gaussLegrendreQuadrature}
\end{equation}

, where $\frac{dx}{d\xi} = \frac{b-a}{2}$ expresses a change of variable in order to yield an integration in $[-1,\,1]$

Then, equation \ref{eq:gaussLegrendreQuadrature} may be approximated by,
	
\begin{equation}
	\int_a^b f(x) dx \approx \frac{a-b}{2} \sum_{i=1}^{n} w_i f\left(\frac{b-a}{2} \xi_i + \frac{a+b}{2}\right)
	\label{eq:gaussLegendreQuadratureApprox}
\end{equation}

In the case of a double integral, it comes that,

\begin{equation}
	\int_{c}^{d} \int_{a}^{b} f\left(x,\,y\right) dx \, dy \approx \frac{(a-b)(d-c)}{4} \sum_{j=1}^{m} \sum_{i=1}^{n} f\left(\frac{b-a}{2} \xi_i + \frac{a+b}{2},\, \frac{d-c}{2} \xi_j + \frac{c+d}{2}\right)
	\label{eq:gaussLegendreDouble}
\end{equation}

In this thesis, a 5-point Gauss Legrendre quadrature was used, where $\xi_i = \{ -0.90618,\, -0.538469,\\ \,0.0,\,0.538469,\,0.90618 \}$ and $w_i = \{ 0.236927,\,0.478629,\,0.568889,\, 0.478629,\,0.236927 \}$.

\section{Boltzman operator}
\label{section:boltzmanOperator}

The boltzmann operator is a smooth, differentiable function that approximates the non-smooth maximum or minimum functions.

\begin{equation}
\mathcal{S}_{\alpha} (x_1,\, \ldots, \,x_n) = \frac{\sum_{i=1}^{n} x_i \exp(\alpha \, x_i)}{\sum_{i=1}^{n} \exp(\alpha \, x_i)}
\label{eq:boltzmannOperator}
\end{equation}

The boltzmann operator has the following properties:

\begin{enumerate}
\item $\lim_{\alpha \to \infty} \mathcal{S}_{\alpha} = \text{max}( x_1,\,\ldots,\,x_n )$
\item $\mathcal{S}_0$ is the arithmetic mean of $( x_1,\,\ldots,\,x_n )$.
\item $\lim_{\alpha \to -\infty} \mathcal{S}_{\alpha} = \text{min}( x_1,\,\ldots,\,x_n )$
\end{enumerate}